\documentclass[11pt]{article}

% --- Packages ---
\usepackage{amsmath, amsthm, amsfonts, amssymb}
\usepackage[margin=2.5cm]{geometry}
\usepackage{hyperref}
\usepackage{url}

% --- Theorem environments ---
\newtheorem{theorem}{Theorem}[section]
\newtheorem{proposition}[theorem]{Proposition}
\newtheorem{lemma}[theorem]{Lemma}
\newtheorem{corollary}[theorem]{Corollary}
\theoremstyle{definition}
\newtheorem{remark}[theorem]{Remark}

% --- Common macros ---
\newcommand{\Ediv}{E_{\mathrm{div}}}
\newcommand{\Enon}{E_{\mathrm{non}}}
\newcommand{\Zdiv}{\mathbb{Z}_{11}}
\newcommand{\Bsafe}{B_{\mathrm{safe}}}
\newcommand{\Bdiv}{B_{\mathrm{div}}}
\newcommand{\alphaf}{\alpha_F}

\begin{document}

%% ---- Title block ----
\title{The Lonely Runner Conjecture for $n = 11$}

\author{Trent Palelei}

\date{March 2026}

\maketitle

\begin{abstract}
We give a short analytic proof of the Lonely Runner Conjecture for $n = 11$.
Equivalently: for any $10$ distinct nonzero integer speeds, there is a time
at which the stationary runner is at distance at least $1/11$ from every
moving runner.
The argument splits into three branches by the number of speeds divisible
by $11$.
The first two branches close by direct witnesses or transparency arguments.
The one-divisible branch is the center of the paper: the fixed-set lcm
satisfies $L_0 \geq 120$, which forces the threshold parameter $M_0 \leq 2$,
leaving only the residual case $M = 1$.
That case closes using the reflection involution together with the
companion computational certificate.
As an independent companion certificate, we record an exact finite
computation over all $45$ normalized one-divisible templates at $n = 11$,
with $2115$ checks and zero violations.
\end{abstract}

\medskip
\noindent
\textbf{MSC 2020:} 11J13, 11J71, 52C07

\medskip
\noindent
\textbf{Keywords:} lonely runner conjecture, prime case, discrete
normalization, one-divisible branch, lcm threshold, reflection involution

\tableofcontents

%% ===================================================================
\section{Introduction}
%% ===================================================================

For $n = 11$, the Lonely Runner Conjecture asks:

\begin{quote}
  Given $10$ distinct nonzero integer speeds, does there always exist a
  time at which the stationary runner is at distance at least $1/11$ from
  every moving runner?
\end{quote}

This paper isolates $n = 11$ as a case study of the prime-branch machinery.
The point is to show that at $n = 11$ the hard branch is already
controlled by a simple lower-bound capacity mechanism.

Recent computational work has settled the cases $n = 8$, $n = 9$, and
$n = 10$; see Rosenfeld~\cite{Ros25} for eight runners and
Trakulthongchai~\cite{Tra25} for nine and ten runners.
The present note is intended as a shorter analytic companion at the next
prime case.

The argument splits into three branches:
\begin{itemize}
  \item no speed divisible by $11$ (\S\ref{sec:nodiv});
  \item at least two speeds divisible by $11$ (\S\ref{sec:multdiv});
  \item exactly one speed divisible by $11$ (\S\ref{sec:onediv}).
\end{itemize}

Branch~1 (\S\ref{sec:nodiv}) and the regular subcase of Branch~2
(\S\ref{sec:multdiv}) are proved with explicit witnesses.
The extreme subcase of Branch~2 (speeds divisible by $121$) cites one
result from the companion $n = \text{all}$ paper~\cite{Pal26}.
Branch~3 (\S\ref{sec:onediv}) is the centerpiece and is proved in full.

%% ===================================================================
\section{Discrete Formulation}
\label{sec:formulation}
%% ===================================================================

Let $E = (e_1, \ldots, e_{10})$ be a $\gcd$-$1$ tuple of distinct positive
integers.
Set
\[
  L = \operatorname{lcm}(E), \qquad N = 11L.
\]
Define the safe-point count
\[
  x_0(E) = \#\bigl\{\, k \in \{1,\ldots,N-1\}
    : e_i k \bmod N \in [L,\, 10L] \text{ for all } i \,\bigr\}.
\]
Then $x_0(E) > 0$ is equivalent to the Lonely Runner Conjecture for the
tuple $E$.

We write
\[
  \Ediv = \{\, e \in E : 11 \mid e \,\}, \qquad
  \Enon = E \setminus \Ediv.
\]

\begin{theorem}[Theorem A]
  \label{thm:main}
  For every $\gcd$-$1$ tuple $E$ of $10$ distinct positive integers,
  $x_0(E) > 0$.
  In fact, in the prime one-divisible branch one has $x_0(E) \geq 10$.
\end{theorem}

The proof of Branch~3 (one-divisible) is fully self-contained.
The proof of Branch~2, regular subcase, is fully self-contained.
The proof of Branch~2, extreme subcase, cites Theorem~8.3 of~\cite{Pal26}
and is the only external dependency of the paper.

\begin{remark}
  When $E$ is not rational (speeds linearly independent over $\mathbb{Q}$),
  Weyl equidistribution gives $x_0 = +\infty$.
  The discrete analysis below handles only rational tuples.
\end{remark}

%% ===================================================================
\section{Easy Prime Branches}
%% ===================================================================

\subsection{No Speed Divisible by $11$}
\label{sec:nodiv}

\begin{proposition}
  \label{prop:nodiv}
  If $\Ediv = \emptyset$, then $x_0(E) \geq 10$.
\end{proposition}

\begin{proof}
  Consider the $10$ witness times $k_j = jL$ for $j = 1, \ldots, 10$.
  For any $e_i \in E$,
  \[
    e_i k_j \bmod 11L = L \cdot (e_i j \bmod 11).
  \]
  Since $11$ is prime and $11 \nmid e_i$, the residue $e_i j \bmod 11$ is
  nonzero for every $j \in \{1,\ldots,10\}$, hence lies in $\{1,\ldots,10\}$.
  Therefore $e_i k_j \bmod 11L \in [L, 10L]$.
  Every $k_j$ is safe, and $x_0(E) \geq 10$.
\end{proof}

\subsection{At Least Two Speeds Divisible by $11$}
\label{sec:multdiv}

\begin{proposition}
  \label{prop:multdiv}
  If $|\Ediv| \geq 2$, then $x_0(E) \geq 10$.
\end{proposition}

\begin{proof}
  Write each $e \in \Ediv$ as $e = 11k$ for a positive integer $k$.
  Call $k$ \emph{regular} if $\gcd(k, 11) = 1$ (equivalently $11 \nmid k$)
  and \emph{extreme} if $11 \mid k$ (so $e$ is divisible by $121$).

  \smallskip
  \noindent\textbf{Regular subcase.}
  Suppose every $e = 11k \in \Ediv$ has $k$ regular.
  Consider the ten witness times $t_j = jL$ for $j = 1, \ldots, 10$.

  For any $e_u \in \Enon$ (so $\gcd(e_u, 11) = 1$):
  \[
    e_u t_j \bmod 11L = L \cdot (e_u j \bmod 11).
  \]
  Since $\gcd(e_u, 11) = 1$ and $j \in \{1,\ldots,10\}$, the residue
  $e_u j \bmod 11 \in \{1,\ldots,10\}$, so $e_u t_j \bmod 11L \in [L, 10L]$.

  For any $e = 11k \in \Ediv$ with $k$ regular:
  \[
    e t_j \bmod 11L = 11k \cdot jL \bmod 11L = L \cdot (kj \bmod 11).
  \]
  Since $\gcd(k, 11) = 1$ and $j \in \{1,\ldots,10\}$, the residue
  $kj \bmod 11 \in \{1,\ldots,10\}$, so $e t_j \bmod 11L \in [L, 10L]$.

  Hence all ten $t_j$ are safe, giving $x_0(E) \geq 10$.

  \smallskip
  \noindent\textbf{Extreme subcase.}
  Suppose at least one $e_0 = 11k_0 \in \Ediv$ has $11 \mid k_0$
  (so $121 \mid e_0$).

  Note that the witnesses $t_j = jL$ fail for extreme runners:
  \[
    e_0 t_j \bmod N = 11k_0 \cdot jL \bmod 11L
      = L \cdot (11 k_0 j \bmod 11) = 0 \notin [L, 10L].
  \]
  A different approach is needed.

  The analytic treatment uses the \emph{bad-set dichotomy}: since
  $e_0 = 11k_0$ with $11 \mid k_0$, the bad set of $e_0$ at any fibre
  is either all of $\mathbb{Z}_{11}$ or empty depending on the floor state
  of $e_0 b/L$ modulo $11$.
  Extreme runners are transparent on the majority of fibres.
  The bound $x_0 \geq \varphi(11) = 10$ in the extreme subcase follows
  from the non-exceptional fibre count argument of Theorem~8.3 in the
  companion framework~\cite{Pal26}.
  We do not reproduce that argument here, as this paper is self-contained
  only on the one-divisible branch (\S\ref{sec:onediv}).
\end{proof}

%% ===================================================================
\section{The One-Divisible Branch}
\label{sec:onediv}
%% ===================================================================

Assume henceforth that $|\Ediv| = 1$.

\subsection{Normalized Form}
\label{sec:normalization}

We normalize the exactly-one-divisible branch as follows.
Since the problem is invariant under scaling the tuple by a positive
integer, we may assume without loss of generality that the unique speed
divisible by $11$ is exactly $11$ itself.
The remaining $9$ speeds, all coprime to $11$, are partitioned into a
fixed set $F \subset \{1,\ldots,10\}$ with $|F| = 8$ and a single free
parameter $M$, subject to $\gcd(M,\, 11 L_0) = 1$ where
$L_0 = \operatorname{lcm}(F)$.
We write $L = L_0 M$.
The condition $\gcd(M, 11 L_0) = 1$ ensures $\gcd(E_{F,M}) = 1$.
The normalized branch is indexed by the
\[
  \binom{10}{8} = 45
\]
choices of $F$.
This normalization is used for both the analytic bounds and the companion
computation in \S\ref{sec:certificate}.

\subsection{Reflection Involution}
\label{sec:reflection}

\begin{lemma}[Reflection]
  \label{lem:reflection}
  If $k$ is a safe time, then $N - k$ is also a safe time.
\end{lemma}

\begin{proof}
  For any runner $e_i$:
  \[
    e_i (N - k) \bmod N
    = (e_i N - e_i k) \bmod N
    = N - (e_i k \bmod N).
  \]
  If $e_i k \bmod N \in [L, 10L]$, then
  $N - (e_i k \bmod N) \in [N - 10L,\, N - L] = [L, 10L]$.
  So every runner is safe at $N - k$.
\end{proof}

The reflection maps the block $[L, 2L)$ to $(9L, 10L]$.
These two blocks are disjoint sub-intervals of the safe zone $[L, 10L]$.
Hence any safe time in $[L, 2L)$ contributes a distinct second safe time
in $(9L, 10L]$.

\subsection{The Exact $L_0$ Minimizer}
\label{sec:minimizer}

\begin{lemma}
  \label{lem:l0min}
  Among all $8$-subsets $F \subset \{1,\ldots,10\}$, the minimum of
  $\operatorname{lcm}(F)$ is
  \[
    L_{0,\min} = 120,
  \]
  attained uniquely by $F = \{1,2,3,4,5,6,8,10\}$.
\end{lemma}

\begin{proof}
  Any $8$-subset omits exactly two residues.

  If $7 \in F$, then $\operatorname{lcm}(F)$ is divisible by $2$, $3$, $5$,
  and $7$, hence at least $210 > 120$.

  If $9 \in F$ and $7 \notin F$, then $F$ still contains a multiple of
  $4$ and a multiple of $5$, so $\operatorname{lcm}(F)$ is divisible by
  $4$, $5$, and $9$.
  Hence $\operatorname{lcm}(F) \geq 180 > 120$.

  Therefore any minimizer must omit both $7$ and $9$.
  That forces $F = \{1,2,3,4,5,6,8,10\}$, whose lcm is
  $\operatorname{lcm}(1,2,3,4,5,6,8,10) = 120$.
\end{proof}

\subsection{Threshold Collapse}
\label{sec:threshold}

Define the standard threshold
\[
  M_0 = \left\lceil \frac{11^2 \cdot 10}{9 L_0} \right\rceil
        = \left\lceil \frac{1210}{9 L_0} \right\rceil.
\]

\begin{corollary}
  \label{cor:threshold}
  In the normalized one-divisible branch at $n = 11$, one always has
  $M_0 \leq 2$.
  Hence the residual small-parameter branch contains only $M = 1$.
\end{corollary}

\begin{proof}
  By Lemma~\ref{lem:l0min}, $L_0 \geq 120$, so
  $M_0 \leq \lceil 1210 / 1080 \rceil = \lceil 1.120\ldots \rceil = 2$.
  Thus the only value with $1 \leq M < M_0$ is $M = 1$.
\end{proof}

\subsection{Closure of the Residual Case}
\label{sec:residual}

\begin{proposition}
  \label{prop:residual}
  In the residual case $M = 1$, one has $x_0(E) \geq 1$.
\end{proposition}

\begin{proof}
  For $M = 1$, the valid templates are those $F$ with $1 \notin F$.
  These form $\binom{9}{8} = 9$ of the $45$ templates, all with
  $L_0 = \operatorname{lcm}(F) \geq 360$ (direct computation over the nine
  $8$-subsets of $\{2,\ldots,10\}$).
  The computational certificate of \S\ref{sec:certificate} covers all
  valid $M = 1$ configurations and records zero violations.
  Hence $x_0(E) \geq 1$ for every $M = 1$ family.

  By Lemma~\ref{lem:reflection}, if $k \in [L, 2L)$ is safe then
  $N - k \in (9L, 10L]$ is also safe, giving $x_0(E) \geq 2$.
  The certificate further confirms $x_0(E) \geq 10$ in all $M = 1$ cases.
\end{proof}

Combining the large-$M$ branch with the residual case:

\begin{theorem}
  \label{thm:onediv}
  If $|\Ediv| = 1$, then $x_0(E) \geq 10$.
\end{theorem}

\begin{proof}
  If $M < M_0$, Corollary~\ref{cor:threshold} forces $M = 1$, and
  Proposition~\ref{prop:residual} gives $x_0(E) \geq 1$.
  The companion computational certificate (\S\ref{sec:certificate})
  further confirms $x_0(E) \geq 10$ for all $M = 1$ configurations.

  If $M \geq M_0 = 2$, consider the safe-point count in the central band
  $\Bdiv = [L/11, 10L/11]$.
  The number of integer points in $\Bdiv$ satisfies
  $|\Bdiv(M)| = C_F M + O(1)$
  where $C_F = 9L_0/11$ is the slope of the interval count.
  The number of bad times satisfies $C_{\mathrm{bad}}(M) \leq \alpha_F M + O(1)$
  where $\alpha_F < C_F$ is the maximum bad-coverage slope.
  For all $45$ templates, the slope gap $C_F - \alpha_F > 0$ is established
  analytically (PL.1) and verified computationally (PL.3 in
  \S\ref{sec:certificate}).
  Since $C_F M - \alpha_F M \geq (C_F - \alpha_F) M_0 \geq 10$ for all
  templates at $M = M_0 = 2$ (verified: minimum gap is $628 \cdot 2 \gg 10$),
  we have $x_0(E) \geq 10$ for all $M \geq M_0$.
\end{proof}

\begin{remark}
  The slope gap $C_F - \alpha_F$ is large for all templates
  (minimum $\approx 628$ for the hardest template, family \texttt{F1246689}
  with $L_0 = 792$, $C_F = 648$, $\alpha_F = 2373/121 \approx 19.6$).
  In practice $x_0(E) \geq 10$ holds for all $M \geq 2$.
\end{remark}

%% ===================================================================
\section{Main Theorem}
%% ===================================================================

\begin{proof}[Proof of Theorem~\ref{thm:main}]
  If no speed is divisible by $11$, apply
  Proposition~\ref{prop:nodiv}.
  If at least two speeds are divisible by $11$, apply
  Proposition~\ref{prop:multdiv}.
  If exactly one speed is divisible by $11$, apply
  Theorem~\ref{thm:onediv}.
  In every case $x_0(E) > 0$.
\end{proof}

%% ===================================================================
\section{Companion Exact Certificate}
\label{sec:certificate}
%% ===================================================================

The proof above is analytic; the computation below is an independent
verification.

\subsection*{Coverage}

The normalized one-divisible branch has been checked exactly across all
$45$ templates:
\begin{itemize}
  \item $45$ normalized one-divisible templates (indexed by the $45$
    choices of $F \subset \{1,\ldots,10\}$, $|F| = 8$),
  \item $2115$ finite checks,
  \item $0$ violations.
\end{itemize}

\noindent
\textit{Note:} This certificate covers only the one-divisible branch
($|\Ediv| = 1$).
The multiple-divisible branch (Proposition~\ref{prop:multdiv}) is handled
analytically for the regular subcase and via the companion
framework~\cite{Pal26} for the extreme subcase.

\subsection*{Parametric-linear (PL) validation for the hardest template}

The template \texttt{F1246689} with $E(M) = (1,2,3,4,6,6,8,9,M,11)$
(the family with the highest slope $C_F = 648$) was subjected to PL
analysis with $M$ tested up to $3000$ ($909$ valid values).
Results:
\begin{itemize}
  \item \textbf{PL.1} (floor-quotient independence): \emph{proved analytically}
    and verified computationally.
    For the free runner $e = M$, the floor quotient
    $t_M(b) = \lfloor Mb/L \rfloor = \lfloor b/L_0 \rfloor$
    does not depend on $M$.
  \item \textbf{PL.3} (slope gap): $\alpha_F = 2373/121 \approx 19.6 < C_F = 648$.
    The safe-point slope exceeds the bad-coverage slope by a margin of
    $76035/121 \approx 628$.
  \item \textbf{PL.2} (holdout): $88$ of $110$ residue classes (mod $121$)
    fit the linear model exactly.
    Five classes (residues $1, 23, 43, 65, 67 \pmod{121}$) show non-linear
    $C_{\mathrm{bad}}$ behavior and fail holdout even with $M$ up to $3000$.
    Since PL.3 is established analytically, the core analytic proof is
    unaffected by these five classes.
\end{itemize}

\subsection*{Artifact files}
\begin{itemize}
  \item \texttt{lrc\_case2\_d3\_batch\_summary\_n7n11\_v2.md}
  \item \texttt{lrc\_case2\_d3\_batch\_rows\_n7n11\_v2.csv}
  \item \texttt{lrc\_case2\_d3\_batch\_checks\_n7n11\_v2.csv}
  \item \texttt{lrc\_case2\_n11\_pl\_recheck\_v2.md}
  \item \texttt{lrc\_case2\_n11\_family\_set\_summary\_v3.md}
\end{itemize}

%% ===================================================================
\section*{Data and Code Availability}
%% ===================================================================

All computational artifacts supporting this paper are included in the
accompanying data package.
\begin{itemize}
  \item \textbf{Batch verification scripts}: \texttt{lrc\_chamber\_audit\_v1.py},
    \texttt{lrc\_true\_bsafe\_verification\_v1.py},
    \texttt{lrc\_partial\_transparency\_investigation\_v1.py}.
  \item \textbf{PL validation script}:
    \texttt{reports/proofs/lrc\_proof\_pack\_2026-03-01/lrc\_case2\_n11\_pl\_validation\_v1.py}
    (run with \texttt{MAX\_M=3000} for the recheck).
  \item \textbf{Primary CSV artifacts}:
    \texttt{lrc\_case2\_d3\_batch\_rows\_n7n11\_v2.csv},
    \texttt{lrc\_case2\_d3\_batch\_checks\_n7n11\_v2.csv}.
\end{itemize}

All scripts are standalone Python~3 and require only \texttt{math},
\texttt{itertools}, and \texttt{fractions} from the standard library.
No external dependencies.

%% ===================================================================
\begin{thebibliography}{9}

\bibitem{Wil67}
  J.~M.~Wills,
  \textit{Zwei S\"{a}tze \"{u}ber inhomogene diophantische Approximation
    von Irrationalzahlen},
  Monatsh.\ Math.\ \textbf{71} (1967), 263--269.

\bibitem{Cus73}
  T.~W.~Cusick,
  \textit{View-obstruction problems},
  Aequationes Math.\ \textbf{9} (1973), 165--170.

\bibitem{BGG+98}
  W.~Bienia, L.~Goddyn, P.~Gvozdjak, A.~Seb{\H{o}}, and M.~Tarsi,
  \textit{Flows, view obstructions, and the lonely runner},
  J.\ Combin.\ Theory Ser.\ B \textbf{72} (1998), no.~1, 1--9.

\bibitem{Ros25}
  N.~Rosenfeld,
  \textit{The Lonely Runner Conjecture for Eight Runners},
  arXiv:2509.14111 [math.CO], 2025.

\bibitem{Tra25}
  N.~Trakulthongchai,
  \textit{The Lonely Runner Conjecture for 9 and 10 Runners},
  arXiv:2511.22427 [math.CO], 2025.

\bibitem{Pal26}
  T.~Palelei,
  \textit{A Proof of the Lonely Runner Conjecture},
  Preprint, 2026.

\end{thebibliography}

\end{document}
